\section{Introduction}\label{sec:introduction}

For a finite set \(X\subset\R^2\), write
\[
 u(X)=\#\bigl\{\{x,y\}\subset X:\ \norm{x-y}_2=1\bigr\},
 \qquad
 u(n)=\max_{\substack{X\subset\R^2\\ |X|=n}}u(X).
\]
Thus pairs are unordered and their endpoints are distinct. We prove a fixed
power saving over the exponent \(4/3\).

\begin{theorem}\label{thm:main}
There are absolute constants \(0<C<\infty\) and \(1\le\beta<4/3\) such that
\(u(n)\le Cn^\beta\) for every nonnegative integer \(n\).
\end{theorem}

\Cref{thm:main} gives a positive answer to the question whether this upper
exponent can be lowered by a fixed amount: equivalently,
\(u(n)=O(n^{4/3-\delta})\) for an absolute \(\delta>0\).

\subsection{The unit-distance problem and earlier methods}

Erd\H{o}s introduced the repeated-distance problem in 1946
\cite{erdos1946}. Rescaling a configuration turns any specified positive
distance into a unit distance, so the problem concerns the largest possible
multiplicity of one distance. His lattice construction gives the lower bound
\(n^{1+c/\log\log n}\) for a positive constant \(c\) and sufficiently large
\(n\) \cite[Theorem~2]{erdos1946}.
Spencer, Szemer\'edi and Trotter established the upper bound
\(u(n)=O(n^{4/3})\) in 1984 \cite{sst1984}.
Sz\'ekely gave a crossing-number proof and a general incidence theorem
for families of curves \cite[Theorems~4 and~8]{Szekely1997}.
The crossing inequality was developed independently by Ajtai,
Chv\'atal, Newborn and Szemer\'edi and by Leighton
\cite{ajtaiEtAl1982,leighton1983}.
Sz\'ekely's incidence theorem applies to points and unit circles using two elementary
properties: two unit circles meet in at most two points, and two points
belong to at most two unit circles.
Later work improved the leading constant
\cite{agostonPalvolgyi2022} and described the structure of configurations
near the extremal incidence scale \cite{katzSilier2026}.
Pach, Raz and Solymosi connected the problem to rigidity and obtained
a logarithmic improvement conditional on a rigidity conjecture
\cite[Conjecture~7 and Theorem~8]{pachRazSolymosi2026}.

The historical unit-distance conjecture of Erd\H{o}s proposed the stronger
bound \(u(n)=n^{1+o(1)}\). An OpenAI construction disproved it by giving
at least \(n^{1+\varepsilon}\) unit distances, for some fixed
\(\varepsilon>0\), along an unbounded sequence of sizes
\cite[Theorem~1.1]{openaiUnitLower2026}. Alon et al.\ gave a simplified
account and explained the construction's number-theoretic ancestry
\cite[Sections~1.2--1.3]{alonEtAl2026}. They credit the class-group
ideas to Michel, Soundararajan, Ellenberg and Venkatesh. The construction
combines these ideas with Golod--Shafarevich
number-field towers having prescribed splitting, including the
tower-cutting method of Hajir, Maire and Ramakrishna
\cite{ellenbergVenkatesh2007,golodShafarevich1964,hajirMaireRamakrishna2021}.
Sawin sharpened the construction to \(u(n)\gg n^{1.014114}\) along an
unbounded sequence of sizes \cite[Theorem~1]{sawin2026}.
These are construction lower bounds. Theorem~\ref{thm:main} bounds every
configuration from above, leaving a gap between the known exponents.

Our proof begins with the same point--circle incidence interpretation.
Random cuttings isolate a piece of the incidence graph with controlled
degrees; their sampled-cell construction belongs to the framework of
Clarkson and Clarkson--Shor \cite{Clarkson1987,ClarksonShor1989}.
The next ingredient is information: two points sampled as neighbors of
one center share little information. Entropy methods already played a
role in the distinct-distance estimates of Tardos and Katz--Tardos
\cite{Tardos2003,KatzTardos}. In particular, the conditional-copy
argument of Katz--Tardos retains the marginal law when resampling
within a class. Here that principle enters a short-history prediction
lemma, proved in \Cref{sec:prediction}, in which the tests are drawn
independently of the vertices and only the criterion for accepting the
outcome of a test at a point may depend on a neighboring center.

The arithmetic input is the normalized product formula
\cite{ProductFormula}: for a nonzero algebraic number, the sum of its
logarithmic absolute values over all places of a number field containing
it, with the normalized local-degree weights, is zero. We use the associated
absolute height of algebraic numbers: the weighted sum of the positive
parts of their logarithmic absolute values. For a unit edge, let \(d\)
be the coordinate \(z=x+iy\) of its point minus that of its center; the
corresponding difference in the coordinate \(w=x-iy\) is then
\(d^{-1}\). This reciprocity is the unit-distance equation itself, and the
proof uses it, not only the two pair-counting properties of unit circles on
which Sz\'ekely's incidence theorem rests: points and lines satisfy those
properties, with one in place of two, and attain the exponent \(4/3\), so no
argument that uses only them can lower it
\cite[Section~1.1]{pachRazSolymosi2026}. For two distinct unit edges with a common endpoint and coordinate differences \(d,d'\),
we have \(d\ne d'\). The product formula applied to \(d/d'\) makes
the sum over all places, with the same weights, of
\(\abs{\log\abs{d/d'}_v}\) twice the height of \(d/d'\)
(\Cref{ht:height}); applied to \((d-d')(d^{-1}-d'^{-1})\), the product
of the differences of the other two endpoints in the two coordinates, it
gives the upper bounds for the integrated scores of pairs of edges
(\Cref{bud:upper,ext:pointwise}). The proof combines these bounds with
prediction and, at its last stage, with algebraic independence and
quadratic extensions.

\subsection{The proof in outline}

Suppose that no fixed power saving exists. Set \(t=n^{1/3}\) along a sequence
of counterexamples, so their ordered unit pairs number \(t^{4-o(1)}\).

\paragraph{Incidence and prediction.}
Random cuttings and degree trimming in \Cref{sec:incidence} give an
unbalanced point--circle incidence graph with controlled degrees.
Maximizing over induced bipartite subgraphs then gives expansion: small vertex sets
capture few edges, and near-agreement of labels across edges forces
concentration on one label on both sides.
For the uniform edge law, sample a center from its marginal and then
two neighbors independently; this gives an edge pair at a center. The small point
side and the bound of two common centers per distinct point pair make
the two point endpoints' mutual information \(o(\log t)\).

Section~\ref{sec:prediction} uses this information bound to resample
a point conditional on its answers to a short history of coordinate
tests. A history can be fixed so that the resampled point preserves
acceptance by a neighboring center at an independent fresh test, up to
a vanishing mean error. Resampling preserves the point marginal,
while the \(t^{o(1)}\) history classes make recovery of the original
point unlikely. This supplies a second point for the geometric comparisons.

\paragraph{Height and the bounded alternative.}
Algebraic specialization preserves the unit equations and distinctness
(\Cref{inc:specialization}), making number-field absolute values
available. In coordinates \(z=x+iy\), \(w=x-iy\), a unit edge has
coordinate differences \(d,d^{-1}\). For an absolute value \(v\) and threshold \(s\), record the threshold bit
\(\1_{\{-\log|d|_v>s\}}\). Let \(f(v,s)\) be the minority-bit
probability and \(S=\int f\), integrating over thresholds with the
normalized weights over places. The integrated disagreement of two
edge bits is twice the height of their coordinate difference ratio.

If \(S\) remained bounded, expansion would transfer height control
from edge coordinate differences to normalized point coordinates. At a fresh complex
embedding, this gives positive probability of bounded, separated
samples. Prediction with quantized coordinates then produces a second,
separated point approximately satisfying the unit equations at three
neighbors of the original point. The determinant for three separated
coordinate differences forces the second point close to the first, a contradiction.
Thus \(S\to\infty\) (\Cref{ht:bounded}).

\paragraph{Comparing point and center pairs.}
Section~\ref{sec:heights} gives each vertex two labels recording
closeness in its two coordinates; their pair is its state. Across an
edge, exactly one label agrees, and the bit specifies which one, except
for integrated error \(O(1)=o(S)\). Separate levels with $f\ge\epsilon$ from those with $f<\epsilon$, where $\epsilon=k^{-1/16}$ for a sample size $k\to\infty$. After discarding levels where $f<\epsilon$ and the consistency error is large, expansion supplies a dominant
coordinate label at each remaining level with $f<\epsilon$. A vertex's exception
mass is the total weighted length of the levels where it lacks that
label.

Sections~\ref{sec:budgets} and~\ref{sec:common} compare pairs of edges
at a point and at a center. Their scores reward certain coordinate
matches and charge bit disagreement. The product formula gives upper
bounds for the integrated scores. On levels with $f\ge\epsilon$, prediction and
the geometry of states agreeing in exactly one coordinate give lower
bounds in terms of bit variation. A slight adjustment of the rewards
near bit balance makes the lower bounds a larger multiple of the bit
variation than the upper bounds. The main terms from levels with $f<\epsilon$ cancel when
the two orientations are added. This forces bit disagreement on levels with $f\ge\epsilon$ in
edge pairs at a center to be \(o(S)\). A second comparison, using two points
resampled from the same history, shows that point exceptions on levels with $f<\epsilon$
retain a positive fraction of \(S\); otherwise expansion would force
\(S=o(S)\). These are \Cref{com:variation,com:rare-positive}.

\paragraph{\texorpdfstring{From levels with $f<\epsilon$ to ratios.}{From levels with f<epsilon to ratios.}}
On the remaining levels with $f<\epsilon$, an edge's bit is determined by its
exceptional endpoints, apart from negligible errors. Subtracting the
two bits in an edge pair at a center cancels the center contribution. Thus for point
endpoints \(p,r\) and center \(q\), the logarithmic absolute values of
\(\lambda_{pr}=(z_p-z_q)/(z_r-z_q)\) are approximately differences
of functions of the absolute-value entry obtained by integrating each point's exception indicator over the retained thresholds, with sign $+1$ for minority bit $1$ and $-1$ for minority bit $0$. These functions cancel in
the alternating product
\[
 \frac{\lambda_{pr}\lambda_{p'r'}}{\lambda_{pr'}\lambda_{p'r}}
\]
around a 4-cycle.

Section~\ref{sec:extraction} uses a pointwise product-formula bound
and expansion to control the tail of the exception mass. It removes
errors in measure weighted by exception mass before selecting a common height
band. The result is a graph with at least \(t^{2-o(1)}\) edges on
two copies of a point set of size at most \(t^{1+o(1)}\), whose
ratio heights are at least \(J/3\), with \(J\to\infty\), and whose
alternating products around 4-cycles all have height \(o(J)\). Pairs on real unit
circles containing at least \(t^{9/10}\) points of the set can
also be removed. This exclusion is needed in the last step.

\paragraph{Differentiation and independent square roots.}
Section~\ref{sec:algebra} turns the separation of heights into an
exact obstruction. In an ultraproduct of the number fields, elements
of height \(o(J)\) form a relatively algebraically closed subfield.
The alternating products around 4-cycles lie in this subfield, while every edge ratio
is transcendental over it. Restricting to plane or curve loci of
minimal total dimension lets us choose a fixed complete bipartite
pattern with positions whose coordinates have no algebraic relations over the field of locus coefficients beyond the locus equations.

A derivation vanishing on the negligible-height subfield can be chosen
nonzero on every ratio in this pattern. The relations for alternating products around 4-cycles split
each ratio's derivative divided by that ratio into a difference of terms
attached to its endpoints. Differentiating
\((z_p-z_r)(w_p-w_r)=2-\lambda_{pr}-\lambda_{pr}^{-1}\) therefore
gives a linear system involving square roots. After one side's
positions and derivative data are adjoined to the field of locus
coefficients, a dimension count leaves a point on the other side
with coordinate transcendence degree equal to the dimension of its locus over that enlarged field. In the function field of that locus, a separate
simple zero for each radicand proves independence of their square
classes. When both loci are the same unit circle the radicands are
squares; this is why that case was excluded earlier.

Quadratic Kummer theory permits independent sign changes of these
square roots. Solving the first few equations expresses the needed
derivative terms using only those roots. Changing the sign of an unused
root then forces an edge ratio's derivative divided by that ratio to vanish,
contradicting its choice. This excludes every sequence approaching the
exponent \(4/3\), establishing a fixed positive gap.

Throughout the proof, asymptotic notation refers to the chosen sequence
\(t\to\infty\). Constants denoted by \(c,C\) may change between occurrences.
All logarithms are natural. Every appeal to a slowly varying parameter
will be made after the quantities it must dominate have been specified.
